\section*{Bab \ref{ch:g9:climate-past-and-present} --- Iklim Bumi, Dahulu dan Sekarang}

\begin{solution}{exo:g9:climate-past-and-present:1}
(a) Sekitar \qty{-19}{\degreeCelsius}. (b) Sekitar \qty{15}{\degreeCelsius}.
(c) $15 - (-19) = 34$: sekitar \qty{34}{\degreeCelsius}. (d) Dua di antara:
karbon dioksida, uap air, metana.
\end{solution}

\begin{solution}{exo:g9:climate-past-and-present:2}
(a) \omterm{def:g1:watching-the-weather:weather}{Cuaca}. (b) \omterm{def:g5:weather-and-climate:climate}{Iklim}. (c) \omterm{def:g5:weather-and-climate:climate}{Iklim} (dan \omterm{def:g9:climate-past-and-present:climate-change}{perubahan iklim}: perubahan selama
puluhan tahun). (d) \omterm{def:g1:watching-the-weather:weather}{Cuaca}.
\end{solution}

\begin{solution}{exo:g9:climate-past-and-present:3}
(a) \qty{19}{cm}, sebelas tahun. (b) $19 \div 11 \approx \qty{1.7}{cm}$.
(c) \omterm{def:g4:clouds-rain-snow:precipitation}{Salju} setiap tahun jatuh di atas \omterm{def:g4:clouds-rain-snow:precipitation}{salju} tahun-tahun sebelumnya,
sehingga makin dalam sebuah lapisan, makin awal \omterm{def:g4:clouds-rain-snow:precipitation}{salju} itu turun.
\end{solution}

\begin{solution}{exo:g9:climate-past-and-present:4}
(a) $22.4 - 21.0 = +\qty{1.4}{\degreeCelsius}$. (b) $20.3 - 21.0 =
\qty{-0.7}{\degreeCelsius}$. (c) \qty{0}{\degreeCelsius}. (d) Bulan itu
lebih dingin daripada rata-ratanya selama 1991--2020.
\end{solution}

\begin{solution}{exo:g9:climate-past-and-present:5}
(a) Tahun 1880-an sampai 1930-an, tahun 1950-an, dan tahun 1960-an.
(b) Sekitar \qty{0.25}{\degreeCelsius} dan \qty{0.80}{\degreeCelsius}.
(c) $0.80 -
0.25 = \qty{0.55}{\degreeCelsius}$. (d) Tidak: dasawarsa-dasawarsa dari
tahun 1880-an sampai 1910-an menjadi sedikit lebih dingin, dan tahun
1950-an sedikit lebih dingin daripada tahun 1940-an; tetapi sejak tahun
1970-an setiap dasawarsa lebih hangat daripada dasawarsa sebelumnya.
\end{solution}

\begin{solution}{exo:g9:climate-past-and-present:6}
(a) $(324.62 - 315.98) \div 10 \approx \qty{0.86}{ppm}$ setahun. (b)
$(427.35 - 401.01) \div 10 \approx \qty{2.63}{ppm}$ setahun. (c) $2.63
\div 0.86 \approx 3.1$: sekitar tiga kali lebih cepat. (d) $(427.35 - 315.98)
\div 315.98 \times 100 \approx 35~\%$.
\end{solution}

\begin{solution}{exo:g9:climate-past-and-present:7}
(a) $8 \times 10^5$ tahun. (b) Sekitar \num{120000} sampai \num{130000}
tahun yang lalu: di sana kurva suhu naik di atas 0. (c) Sekitar
\qty{275}{ppm}. (d) Ya: puncak dan \omterm{def:g1:landscapes-around-us:landform}{lembah} kedua kurva muncul pada waktu
yang sama.
\end{solution}

\begin{solution}{exo:g9:climate-past-and-present:8}
(a) $3.2 \times 10 = \qty{32}{mm} = \qty{3.2}{cm}$. (b) $2100 - 2025 =
75$ tahun; $75 \times 3.2 = \qty{240}{mm}$, sekitar \qty{24}{cm} (dan
lebih jika kenaikannya makin cepat). (c) Air \omterm{def:g2:water-on-earth:ocean}{laut} memuai ketika
menghangat; es di daratan mencair dan airnya mengalir ke \omterm{def:g2:water-on-earth:ocean}{laut}. (d) Tidak:
\omterm{def:g8:ice-on-earth:sea-ice}{es laut} sudah mengapung di \omterm{def:g2:water-on-earth:ocean}{laut}, sehingga mencairnya tidak menaikkan
\omterm{def:g4:world-physical-map:altitude}{permukaan laut}.
\end{solution}

\begin{solution}{exo:g9:climate-past-and-present:9}
(a) 90. (b) 66. (c) 99. (d) \emph{Mungkin} berarti sekurang-kurangnya 66
peluang dari 100, bukan kepastian: surat kabar itu membuat pernyataan itu
terdengar lebih pasti daripada yang dikatakan para ilmuwan.
\end{solution}

\begin{solution}{exo:g9:climate-past-and-present:10}
Satu musim dingin yang dingin di satu tempat adalah \emph{\omterm{def:g1:watching-the-weather:weather}{cuaca}}: musim
itu menunjukkan apa yang terjadi selama beberapa bulan di satu titik.
\omterm{def:g9:climate-past-and-present:global-warming}{Pemanasan global} adalah perubahan \emph{\omterm{def:g5:weather-and-climate:climate}{iklim}}, yang diukur selama
puluhan tahun, dan bersifat \emph{global}: pemanasan itu adalah rata-rata
di seluruh Bumi. Musim dingin yang dingin di sini dapat terjadi sementara
Bumi secara keseluruhan lebih hangat daripada sebelumnya, dan rata-rata
per dasawarsa terus naik.
\end{solution}

\begin{solution}{exo:g9:climate-past-and-present:11}
(a) $280 - 190 = \qty{90}{ppm}$, dalam sekitar $2 \times 10^4$ tahun. (b) $427
- 280 = \qty{147}{ppm}$, dalam $2025 - 1750 = 275$ tahun. (c) Kenaikan
kedua lebih besar (147 dibandingkan dengan 90~ppm). Lajunya: $90 \div (2 \times 10^4) =
4.5 \times 10^{-3}$~ppm setahun, dibandingkan dengan $147 \div 275 \approx 5.3 \times
10^{-1}$~ppm setahun: sekitar 120 kali lebih cepat. (d) Bahwa pada tahun
2019 kadarnya lebih tinggi daripada kapan pun dalam sekurang-kurangnya dua
juta tahun (keyakinan tinggi).
\end{solution}

\begin{solution}{exo:g9:climate-past-and-present:12}
(a) Keduanya hampir sama: 1.07 dibandingkan dengan
\qty{1.06}{\degreeCelsius}. Perkiraan terbaik itu menyatakan bahwa
pengaruh manusia menjelaskan kira-kira seluruh pemanasan yang teramati.
(b) Tidak: penyebab-penyebab alami bersama-sama mengubah suhu paling
banyak \qty{0.1}{\degreeCelsius}, sepuluh kali lebih kecil daripada
pemanasan yang teramati. (c) Sekurang-kurangnya 66 peluang dari 100 bahwa
pemanasan akibat manusia berada di antara 0.8 dan
\qty{1.3}{\degreeCelsius}.
\end{solution}

\begin{solution}{pb:g9:climate-past-and-present:1}
\textbf{1.} Karena periode itu adalah periode acuan: anomali diukur dari
rata-ratanya sendiri.

\textbf{2.} Lebih dingin, sebesar \qty{0.22}{\degreeCelsius}.

\textbf{3.} Lebih hangat, sebesar \qty{0.98}{\degreeCelsius}.

\textbf{4.} $0.98 - (-0.22) = \qty{1.20}{\degreeCelsius}$.

\textbf{5.} $427 - 316 = \qty{111}{ppm}$.

\textbf{6.} $111 \div 316 \times 100 \approx 35~\%$.

\textbf{7.} $427 - 300 = \qty{127}{ppm}$.

\textbf{8.} Karbon dioksida menyerap sebagian cahaya inframerah yang
dipancarkan permukaan yang hangat dan mengirimkan sebagiannya kembali ke
bawah. Makin banyak karbon dioksida di udara, makin banyak energi yang
kembali ke permukaan, sehingga permukaan menghangat.

\textbf{9.} $2025 - 1993 = 32$ tahun; $32 \times 3.2 \approx
\qty{102}{mm}$, sekitar \qty{10}{cm}.

\textbf{10.} Air \omterm{def:g2:water-on-earth:ocean}{laut} yang menghangat memuai; es di daratan (\omterm{def:g8:ice-on-earth:glacier}{gletser},
\omterm{def:g8:ice-on-earth:ice-sheet}{lapisan es}) mencair ke \omterm{def:g2:water-on-earth:ocean}{laut}.

\textbf{11.} Ya, besarnya sama: sekitar \qty{1.2}{\degreeCelsius}
dibandingkan dengan \qty{1.09}{\degreeCelsius}. Keduanya sedikit berbeda
karena periodenya tidak sama (1880--1899 sebagai ganti 1850--1900;
2015--2024 sebagai ganti 2011--2020: tahun-tahun yang lebih baru lebih
hangat), dan karena IPCC menggabungkan beberapa catatan, tidak hanya
catatan NASA.

\textbf{12.} ``Sejak akhir abad ke-19, permukaan Bumi telah menghangat
sekitar \textbf{\qty{1.2}{\degreeCelsius}}.''
\end{solution}
